% examples/tikz-demo.tex
%
% Exercise the Kyoto TikZ helper library (tikzlibrarykyoto.code.tex).
% Build from the repository root with:
%
%   make
%
% which runs latexmk (LuaLaTeX) via examples/latexmkrc for both example
% PDFs -- no need to install anything into TEXMF or set TEXINPUTS by
% hand. Equivalently, from this directory: `latexmk tikz-demo.tex`.
%
\documentclass[aspectratio=169]{beamer}
\usetheme{Kyoto}
\usepackage{tikz}
\usetikzlibrary{kyoto}
\usetikzlibrary{automata}

\begin{document}

% 1. Boxes and a straight arrow ---------------------------------------
\begin{frame}{Boxes and a straight arrow}
  \centering
  \begin{tikzpicture}
    \node[kyoto/box] (a) {Learner};
    \node[kyoto/box, right=3cm of a] (b) {Oracle};
    \draw[kyoto/arrow] (a) -- (b);
  \end{tikzpicture}
\end{frame}

% 2. Curved arrows and labels ------------------------------------------
\begin{frame}{Curved arrows and labels}
  \centering
  \begin{tikzpicture}[node distance=2cm]
    \node[kyoto/box] (learner) {Learner};
    \node[kyoto/box, right=of learner] (oracle) {Oracle};

    \draw[kyoto/arrow] (learner) to[bend left]
      node[kyoto/label, above] {Query} (oracle);
    \draw[kyoto/arrow alert] (oracle) to[bend left]
      node[kyoto/label, below] {Answer} (learner);
  \end{tikzpicture}
\end{frame}

% 3. Box variants: normal vs alert vs secondary -------------------------
\begin{frame}{Box variants}
  \centering
  \begin{tikzpicture}
    \node[kyoto/box] (n) {Normal};
    \node[kyoto/box alert, right=1.5cm of n] (al) {Alert};
    \node[kyoto/box secondary, right=1.5cm of al] (se) {Secondary};
  \end{tikzpicture}
\end{frame}

% 3b. Filled (default) vs outline box styles, side by side ---------------
\begin{frame}{Box styles: filled (default) vs outline}
  \centering
  \begin{tikzpicture}[node distance=1.2cm]
    \node[kyoto/box] (n) {Normal};
    \node[kyoto/box alert, right=1.2cm of n] (al) {Alert};
    \node[kyoto/box secondary, right=1.2cm of al] (se) {Secondary};

    \node[kyoto/box outline, below=of n] (on) {Outline};
    \node[kyoto/box outline alert, right=1.2cm of on] (oal) {Outline alert};
    \node[kyoto/box outline secondary, right=1.2cm of oal] (ose) {Outline secondary};
  \end{tikzpicture}
\end{frame}

% 3c. Outline boxes in a realistic two-node diagram -----------------------
\begin{frame}{Outline boxes in a diagram}
  \centering
  \begin{tikzpicture}[node distance=3cm]
    \node[kyoto/box outline] (learner) {Learner};
    \node[kyoto/box outline alert, right=of learner] (oracle) {Oracle};
    \draw[kyoto/arrow] (learner) -- (oracle);
  \end{tikzpicture}
\end{frame}

% 4. Arrow variants: normal vs alert vs secondary, plus dashed ----------
\begin{frame}{Arrow variants}
  \centering
  \begin{tikzpicture}
    \node[kyoto/box] (a1) {A}; \node[kyoto/box, right=2.5cm of a1] (b1) {B};
    \node[kyoto/box, below=1cm of a1] (a2) {A}; \node[kyoto/box, right=2.5cm of a2] (b2) {B};
    \node[kyoto/box, below=1cm of a2] (a3) {A}; \node[kyoto/box, right=2.5cm of a3] (b3) {B};
    \node[kyoto/box, below=1cm of a3] (a4) {A}; \node[kyoto/box, right=2.5cm of a4] (b4) {B};

    \draw[kyoto/arrow] (a1) -- (b1);
    \draw[kyoto/arrow alert] (a2) -- (b2);
    \draw[kyoto/arrow secondary] (a3) -- (b3);
    \draw[kyoto/arrow dashed] (a4) -- (b4)
      node[kyoto/label, midway, above] {optional};
  \end{tikzpicture}
\end{frame}

% 5. Positioning library usage (below=of, right=of together) -----------
\begin{frame}{Positioning}
  \centering
  \begin{tikzpicture}[node distance=1.5cm]
    \node[kyoto/box] (top) {Top};
    \node[kyoto/box, below=of top] (mid) {Middle};
    \node[kyoto/box secondary, right=2cm of mid] (side) {Side};

    \draw[kyoto/arrow] (top) -- (mid);
    \draw[kyoto/arrow secondary] (mid) -- (side);
  \end{tikzpicture}
\end{frame}

% 6. Callouts: pointer geometry stays ordinary TikZ ----------------------
\begin{frame}{Callouts}
  \centering
  \begin{tikzpicture}[node distance=2.5cm]
    \node[kyoto/box] (learner) {Learner};
    \node[kyoto/box, right=of learner] (oracle) {Oracle};
    \draw[kyoto/arrow] (learner) -- (oracle);

    \node[
      kyoto/callout,
      callout absolute pointer={(learner.north)},
      above=1.5cm of learner,
      text width=3cm,
    ] {Proposes a hypothesis};

    \node[
      kyoto/callout alert,
      callout absolute pointer={(oracle.north)},
      above=1.5cm of oracle,
      text width=3cm,
    ] {May return a counterexample};
  \end{tikzpicture}
\end{frame}

% 6b. Callout styles: filled (default) vs outline -------------------------
\begin{frame}{Callout styles: filled (default) vs outline}
  \centering
  \begin{tikzpicture}[node distance=2.5cm]
    \node[kyoto/box] (t1) {Target};
    \node[
      kyoto/callout,
      callout absolute pointer={(t1.south)},
      above=1.7cm of t1,
      text width=2.8cm,
    ] {Filled (default)};

    \node[kyoto/box, right=3.5cm of t1] (t2) {Target};
    \node[
      kyoto/callout outline,
      callout absolute pointer={(t2.south)},
      above=1.7cm of t2,
      text width=2.8cm,
    ] {Outline};
  \end{tikzpicture}
\end{frame}

% 7. Standard TikZ automata + a Kyoto callout ----------------------------
\begin{frame}{Automata (standard TikZ) with a Kyoto callout}
  \centering
  \begin{tikzpicture}[
      ->,
      >={Stealth[length=2.5mm, width=2mm]},
      node distance=3cm,
      initial text={},
      every state/.style={draw=KyotoTheme, text=KyotoFg, fill=KyotoBg},
    ]
    \node[state, initial] (q0) {$q_0$};
    \node[state, accepting, right=of q0] (q1) {$q_1$};

    \draw (q0) edge[bend left] node[kyoto/label, above] {$a$} (q1)
          (q1) edge[bend left] node[kyoto/label, below] {$b$} (q0);

    \node[
      kyoto/callout,
      callout absolute pointer={(q1.north)},
      above right=1.2cm and -0.5cm of q1,
      text width=3cm,
    ] {Accepting state};
  \end{tikzpicture}
\end{frame}

\end{document}
